\subsection*{How to Read This Atlas}


Every theoretical framework is a perspectival being (DoPI, Theorem 9 applied to formalisms). Each entry maps:


\begin{itemize}[nosep,leftmargin=*]
  \item \textbf{SEES:} What the framework is maximally sensitive to — its coherence dimensions
  \item \textbf{NULL SPACE:} What it structurally cannot access — not limitations of current knowledge but architectural exclusions
  \item \textbf{COMPLEMENTS:} Which other frameworks cover its blind spots
  \item \textbf{BOUNDARY:} Where the framework transitions from reliable to unreliable — the edge of its validity
  \item \textbf{NAVIGATIONAL IMPLICATION:} What a navigator using this framework should do about its blind spots — when to switch frameworks, when to accept the limitation, and how this framework's constraints concentrate the remaining information (see the Guide, Parts I-II)
\end{itemize}


The symbol $\emptyset$ marks absolute null spaces — distinctions the framework cannot access in principle, not merely in practice. The symbol $\bullet$ marks partial null spaces — accessible in principle but poorly resolved.


\textbf{Connection to the Guide:} Each framework is a bottleneck (Guide §1.2). Its null space is what the bottleneck excludes. Its SEES entries are the coherence dimensions it admits. The NAVIGATIONAL IMPLICATION makes explicit what the Guide describes abstractly: the oscillation between frameworks IS the expansion-contraction rhythm of consciousness navigating through mathematical configuration space.


\begin{quote}
\itshape
\textit{See also: Doctrine Theorem 19 (Observational Null Space) for the formal proof that every framework has a structural null space; Guide Part V for practical methods of null-space navigation; Ecology §5 Principles 1-4 for the theory of attention that explains why null spaces exist.}
\end{quote}


\bigskip\noindent\makebox[\textwidth]{\color{rulegray}\rule{5cm}{0.3pt}}\bigskip
